\cotefile{10}
\scanpage{139}
Motifs à coefficients sur un corps de nbs.

Soit d'abord $\unsure{\mathcal C|C|\underline C}$ une $\otimes$-catégorie sur un corps $K$. Si $K'$ est
une extension \inserted{\uline{finie}} de $K$, on désigne par $\mathcal C^{K'}$ la catégorie des
objets de $\mathcal C$ sur lesquels $K'$ opère \unsure{linéairement}, \unsure{i.e.}
des $X\in\operatorname{Ob}\mathcal C$ munis d'un hom de $K$-algèbres
$K'\longrightarrow\operatorname{End}(X)$. C'est encore une $\otimes$-catégorie, en définissant
$X\otimes_{K'}Y$ comme le quotient de $X\otimes_KY$ par les
$\lambda\otimes\mathrm{id}_Y-\mathrm{id}_X\otimes\lambda$ $(\lambda\in K')$, et $\underline{\mathrm{Hom}}_{K'}(X,Y)$ comme le
noyau \unsure{dans|des|de} $\underline{\mathrm{Hom}}_K(X,Y)$ des $\underline{\mathrm{Hom}}(\lambda_X,\mathrm{id}_Y)-\underline{\mathrm{Hom}}(\mathrm{id}_X,\lambda_Y)$.

Quand $\mathcal C$ est défini par un schéma en groupes
affine $G$ sur $K$\inserted{, grâce au foncteur fibre $F$ sur $\mathcal C$}, alors $\mathcal C^{K'}$ est défini par
le groupe affine $G_{K'}=G\otimes_K K'$, grâce au foncteur
fibre \inserted{sur $\mathcal C^{K'}$} $F_{K'}$ déduit trivialement de $F$.

On a un $\otimes$-foncteur naturel $\mathcal C\longrightarrow\mathcal C^{K'}$, \unsure{noté|notant}
$X\leadsto X\otimes_K K'$, et dont le composé avec le foncteur
``oubli'' est bien $X\otimes_K\underset{\sim}{K'}$ [où $\underset{\sim}{K'}$ est \deleted{\illegible{1}} \unsure{le'objet|l'objet}
de $\mathcal C$ défini par le \uline{vectoriel} $K'$ sur $K$, -- puisqu'on
a un $\otimes$-foncteur évident Vectoriels \inserted{de dim finie} sur $K\longrightarrow\mathcal C$].

Supposons $\mathcal C$ la catégorie des \inserted{(iso)}motifs ($1/2$ simples) sur un
corps $k$, \unsure{donc} $K=\mathrm Q$, $K'$ un corps de nbs.
Soit $E$ un motif à coefficients dans $K'$, donc
un motif ordinaire, avec $K'\longrightarrow\operatorname{End}_{\mathrm Q}(E_0)$, $E_0$ le $\mathrm Q$-motif \unsure{ss-jacent}.
Alors $\unsure{\check E|\widetilde E}$ \inserted{dans $\mathcal C^{K'}$} est défini par $\widetilde E$ dans $\mathcal C$, avec l'hom
% archive sheet number: 138
